\section{The Serre torus, CM types, and the limits of the class-field slogan}
\label{sec:part-d}

\providecommand{\Res}{\operatorname{Res}}
\providecommand{\MT}{\operatorname{MT}}
\providecommand{\Cl}{\operatorname{Cl}}
\providecommand{\CM}{\operatorname{CM}}
\providecommand{\Tors}{\operatorname{Tors}}
\providecommand{\Isom}{\operatorname{Isom}}
\providecommand{\cor}{\operatorname{cor}}
\providecommand{\res}{\operatorname{res}}
\providecommand{\ord}{\operatorname{ord}}

The distinctions between a finite-modulus Serre group, the connected Serre
protorus, and the Taniyama group are essential here.  We use $S^K$ for the
standard connected Serre torus attached to a Galois CM field $K$, $S$ for its
projective limit, and $S_{\mathfrak m}$ for Serre's possibly disconnected
finite-modulus group.  The latter is a commutative group; the Taniyama group
is generally not.  Statements about the wording of 15:157 and 15:159 below
refer to the supplied AI transcriptions.  The claimed earlier visual check
of ``m\'etab\'eliens'' cannot be independently repeated without the scans.

\subsection{D.1: the integral model and the actual direction of reduction}
\label{subsec:d1}

\subsubsection{The schematic quotient and the narrow class group}

Let $k$ be a number field, $A=\mathcal O_k$, $E^+=A^{\times,+}$, and
$\mathcal T=\Res_{A/\Z}\Gm$.  Let $\mathcal R$ be the schematic closure
of the sections $E^+\subset\mathcal T(\Z)$.  This is the meaning of
``adh\'erence de l'union des sections'' in 15:159.  It is not the profinite
closure occurring in the adjacent class-field-theoretic formula.

\begin{proposition}[Proved: integral finite-modulus Serre group]
\label{prop:d-integral-serre}
The fppf quotient $\mathcal G^0=\mathcal T/\mathcal R$ is a separated,
flat group scheme of finite type over $\Z$.  Its geometric fibres are
connected.  There is an extension of group schemes
\[
  1\longrightarrow\mathcal G^0\longrightarrow\mathcal G
   \longrightarrow\underline{\Cl^+(k)}\longrightarrow1
\]
whose generic fibre is Serre's $S_{\mathfrak m}$ for
$\mathfrak m=\prod_{v\text{ real}}v$.  In particular
$(\mathcal G^0)_\Q=T_{\mathfrak m}
=\Res_{k/\Q}\Gm/\overline{E^+}^{\,\mathrm{Zar}}$.
The assertion does not say that these integral models are tori at ramified
primes or that their representations comprise all CM motives over $k$.
\end{proposition}

\begin{proof}
The coordinate ring of the closure embeds in the product of copies of
$\Z$ obtained by evaluation at the sections.  It is therefore
$\Z$-torsion-free and flat.  Equivalently, $\mathcal R$ is the schematic
closure of the closed generic-fibre subgroup
$\overline{E^+}^{\,\mathrm{Zar}}$.  Multiplication and inversion preserve
this closure: equality on the generic fibre extends over $\Z$, since the
source schemes are flat.  Thus $\mathcal R$ is a flat closed subgroup of
$\mathcal T$.  Representability and finite type of the fppf quotient follow
from \citet[Theorem 4.C, p.~53]{DAnantharaman1973}; separatedness follows
from closedness of $\mathcal R$.  The map $\mathcal T\to\mathcal G^0$ is
faithfully flat, and flatness descends along it.  For an algebraically closed
residue field $\Omega$, the unit group of $A\otimes\Omega$ is a product
of groups of the form $\Gm\times\mathbb A^n$ as varieties, hence is
connected.  Its quotient $\mathcal G^0_\Omega$ is connected as well.

Here is an integral construction of the class-group extension, including
its multiplication.  For every $c\in\Cl^+(k)$ choose an invertible
fractional ideal $\mathfrak a_c$, with $\mathfrak a_0=A$.  Put
\[
 \mathcal P_c=\Isom_{A}(\mathfrak a_c,A),\qquad
 \mathcal Q_c=\mathcal P_c/\mathcal R.
\]
The first is a $\mathcal T$-torsor and the second a
$\mathcal G^0$-torsor; representability follows either from the same
quotient theorem or from fppf descent.  Choose totally positive
$t_{c,d}\in k^\times$ with
$\mathfrak a_c\mathfrak a_d=t_{c,d}\mathfrak a_{c+d}$.  The map
\[
 (u,v)\longmapsto (u\otimes v)\circ[t_{c,d}]
\]
on ideal isomorphisms defines
$\mathcal Q_c\times\mathcal Q_d\to\mathcal Q_{c+d}$.  Any two choices
of $t_{c,d}$ differ by $E^+$, so the map is independent of the choice.
The ratios occurring in the associativity identity are also positive
units; the identity and inverse maps are constructed in the same manner.
Consequently $\mathcal G=\coprod_c\mathcal Q_c$ is the required group
scheme.  Changing the representatives gives the canonically isomorphic
construction by the same ideal-isomorphism torsors.

Over $\Q$ all these torsors are trivial.  Choosing their trivializations
gives precisely the pushout cocycle of
$k^\times/E^+\to I_{\mathfrak m}\to\Cl^+(k)$ under
$k^\times/E^+\to T_{\mathfrak m}(\Q)$.  This is Serre's definition:
\citet[Chapter II, \S1.3 and \S2.2, pp.~II-3--II-9]{DSerre1968}.
\end{proof}

\paragraph{Verification of R4 and the finite-modulus qualification.}
The finite-idelic class-field exact sequence is
\[
 1\longrightarrow\widehat A^\times/\overline{E^+}
 \longrightarrow\Gal(k^{\mathrm{ab}}/k)
 \longrightarrow\Cl^+(k)\longrightarrow1.
\]
Indeed the valuation map on finite ideles has kernel
$\widehat A^\times$, and quotienting by principal ideles with positive
archimedean components gives exactly the narrow ideal class group.
Using all principal ideles instead leaves the archimedean sign group in
the kernel and produces $\Cl(k)$.  This proves the correction asserted in
R4.  Here the bar denotes topological closure; in
Proposition~\ref{prop:d-integral-serre} it denotes schematic closure.

The finite-modulus torus need only be isogenous to the standard connected
Serre torus.  For example, when $k=\Q(i)$ and $\mathfrak m$ has empty
finite support,
\[
 T_{\mathfrak m}=\Res_{k/\Q}\Gm/\mu_4,
 \qquad S^k=\Res_{k/\Q}\Gm.
\]
Their character lattices, with complex conjugation swapping the two
coordinates, are respectively
\[
 \{(a,b)\in\Z^2:a-b\equiv0\pmod4\},\qquad\Z^2.
\]
The first has basis $(1,1),(2,-2)$ and is the direct sum of the trivial
and sign lattices.  In the second, the sum of the invariant and
anti-invariant sublattices has index two.  Hence the lattices, and the
tori, are not isomorphic.  Allowing a sufficiently small congruence
subgroup of the units removes this finite discrepancy.  A fixed modulus
also imposes a fixed conductor condition; taking only the real places
does not include CM motives ramified at finite primes.

\subsubsection{Frobenius, reflex norms, and two false literal assertions}

\begin{proposition}[Refuted literal Frobenius formula; proved correction]
\label{prop:d-frobenius-correction}
The CM type and a uniformizer do not, by the reflex norm alone,
determine the Frobenius element of a CM motive.  They determine its
algebraic Hecke-character part and its principal ideal; the finite-order
character data are also required to determine the element.
\end{proposition}

\begin{proof}
The precise Shimura--Taniyama assertion is an ideal or an idelic
reciprocity assertion.  If $B$ has CM by $E$ of type $\Phi$, all this
endomorphism structure is defined over $k$, and $B$ has good reduction at
$v$, then the Frobenius $\pi_v$ belongs to $E$, its principal ideal is
the appropriate reflex norm of $N_{k/E^*}(v)$, and its slopes satisfy
\begin{equation}
 \frac{\ord_w(\pi_v)}{\ord_w(q_v)}
 =\frac{\#\{\phi\in\Phi_p:\phi\text{ induces }w\}}
 {[E_w:\Q_p]}.
 \label{eq:d-st-slopes}
\end{equation}
Equivalently, the actual Frobenius is the value of the algebraic Hecke
character of $B$ on a uniformizer idele.  The finite-order part of this
character is necessary.  The reflex norm specifies its algebraic part;
it does not by itself remove the unit or root-of-unity ambiguity.
See \citet[Theorem 8.1, Remark 8.6(b), and Theorems 9.3 and 9.10]{DMilneCM2020}
and the particularly explicit
\citet[Theorem 5.1(ii)--(v)]{DCCO2008}.
Here $\pi_v$ is the $q_v$-power Frobenius endomorphism, whose
eigenvalues are those of geometric Frobenius on $H^1$.  The cited
Hecke-character formula uses arithmetic reciprocity on the Tate
module, as specified in \S1.7 of that source.  Dualizing to
cohomology accounts for this convention without changing the
displayed eigenvalues or slopes.

An explicit check is provided by the two curves over $\Q(i)$
\[
 B:y^2=x^3-x,
 \qquad B^{(2)}:y^2=x^3-4x.
\]
They have the same CM type and endomorphism ring $\Z[i]$, and both have
good reduction at the primes above $5$.  Their discriminants are $64$
and $4096$, so good reduction here follows directly from the equations.
Counting points gives
$\#B(\F_5)=8$ and $\#B^{(2)}(\F_5)=4$, and hence Frobenius traces
$-2$ and $2$.  Their Frobenius eigenvalues are
$-1\pm2i$ and $1\pm2i$.  Under the identification of endomorphisms
induced by the quadratic twisting isomorphism, Frobenius changes sign.
Thus the CM type and the chosen prime give the same reflex-norm ideal
but not the same Frobenius element.  The exact enumeration is attached
in \texttt{part\_d/check\_part\_d.py} and its recorded output.
\end{proof}

\begin{proposition}[Refuted: special-fibre gerbe bridge]
\label{prop:d-special-fibre-false}
The proposed morphism from the finite-field motivic gerbe to the
special fibre of the integral group in Lead D1 is not a morphism of
gerbes over the coefficient field, and good reduction does not make
the Frobenius an integral point of that group.
\end{proposition}

\begin{proof}
There is an even simpler obstruction to the proposed special-fibre
bridge in Lead D1.  For $k=\Q$ the integral group above is $\Gm$.
The Tate motive $\Q(-1)$ has good reduction everywhere and Frobenius
eigenvalue $p$ at $p$.  This is not a $\Z_p$-point of $\Gm$, and its
reduction is not a point of $\Gm_{\F_p}$.  Reversing the Frobenius
convention replaces $p$ by $p^{-1}$, which is not integral.  Moreover,
the motivic gerbe of a finite field is a gerbe over the coefficient
field $\Q$, whereas $\mathcal G_{\F_p}$ is a group scheme over
$\F_p$.  These are not two bands over a common base to which the
asserted morphism applies.  Lead D1 in its special-fibre formulation
is therefore false.
\end{proof}

\begin{proposition}[Proved with the stated CM Hodge input: reduction of gerbes]
\label{prop:d-cm-gerbe-reduction}
Reduction of CM motives induces a morphism of gerbes over $\Q$
from the abelian motivic gerbe over $\overline\F_p$ to the CM
motivic gerbe over $\overline\Q$, bound by $P\to S$.  It satisfies
the Langlands--Rapoport local conditions.  These local conditions
alone do not canonically specify its global comparison with a
separately chosen pseudomotivic morphism.
\end{proposition}

\begin{proof}
Fix an extension of $v$ to $\overline\Q$.  Reduction of CM abelian
motives gives an exact tensor functor over $\Q$
\[
 R:\CM(\overline\Q)\longrightarrow
       \operatorname{Mot}_{\mathrm{ab}}(\overline\F_p),
\]
and therefore, by precomposition on fibre functors, a morphism
\begin{equation}
 R^\vee:\operatorname{Mot}_{\mathrm{ab}}(\overline\F_p)^\vee
    \longrightarrow\CM(\overline\Q)^\vee,
 \qquad\text{with band map }\rho:P\longrightarrow S.
 \label{eq:d-reduction-gerbe}
\end{equation}
This is the direction that exists.  Both sides remain over $\Q$.
The construction and its local properties are proved in
\citet[\S7, pp.~39--40]{AMilne2003}.  The Hodge-conjecture hypothesis
on CM abelian varieties used there is the conclusion of
\citet[Theorem I]{AvanHoftenSempliner2026}; this is an explicitly
identified 2026 preprint input, rather than an assumption about general
Hodge classes.  It permits algebraic specialization of the
correspondences and supplies the required abelian Tate and positivity
statements through the implications cited in that source.

The map \eqref{eq:d-reduction-gerbe} satisfies the Langlands--Rapoport
local conditions: the prime-to-$p$ comparison torsors are trivial,
the $p$-adic map is the isocrystal map, and the real map is the
polarized real realization.  These are the statements immediately
preceding \citet[Theorem 7.2]{AMilne2003}.  They do not uniquely
identify a global comparison morphism.  That theorem says that the
classes of compatible $(P,S)$-gerbes form a principal homogeneous
space under $\varprojlim_K\Cl(K)$; Remark 7.3 gives equality at each
finite CM level, and Remark 7.4 warns that the chosen local objects add
further comparison data.  Thus the motivic reduction is a realization
of the Langlands--Rapoport type of morphism.  Identifying it with a
particular separately chosen pseudomotivic morphism requires the
corresponding global comparison choice.  The local statements alone
cannot prove that stronger assertion.
\end{proof}

\paragraph{Weil numbers: the field of reduction matters.}
At the level of Weil-number germs, Honda--Tate and CM lifting after
a finite residue-field extension and isogeny already suffice.  This
is precisely the formulation of \citet[Theorem (1.7)]{DOort2009};
its Theorem (1.9) recalls the Honda--Tate classification.  For exact
eigenvalues over a fixed finite field one must distinguish lifting
the simple isogeny class itself from obtaining its eigenvalue with
larger multiplicity.  The first can fail, whereas the second always
works when the number field and CM type may vary.  We prove this
stronger and precise statement after explaining the obstruction.

For an isotypic $B_0/\F_q$ with a specified CM field
$E\hookrightarrow\End^0(B_0)$ of degree $2\dim B_0$, and a
$p$-adic CM type $\Phi_p$, the \emph{residual reflex condition} is
\eqref{eq:d-st-slopes} together with the requirement that the residue
field of the reflex field at the specified $p$-adic place embed in
$\F_q$.  The existence, up to an $E$-linear $\F_q$-isogeny, of a CM
lifting over a normal characteristic-zero local domain with exactly
that residue field is equivalent to this condition:
\citet[Theorem (4.4)]{DOort2009}.  More precisely,
\citet[Theorem 1.6 and Proposition 5.4]{DCCO2008} produce an
abelian variety over a number field, with good reduction and exactly
the prescribed residue field.  This is a finite criterion once the
CM structure is specified: enumerate the CM types, compute their
stabilizers and local reflex fields, and test the displayed slope
identities and residue degrees.  The squarefree-Frobenius version over
a number field is \citet[Definition 2.14 and Theorem 2.16]{DBergstromEtAl2023},
which cites Chai--Conrad--Oort, Theorems 2.1.6 and 2.5.3.  The latter
monograph's full text was not obtained; the corresponding general
theorem and proof were opened in the authors' 26 July 2008 version,
Theorem 1.6 and Proposition 5.4, pp.~6 and 31--33.  The author-uploaded
copy has erroneous hosting-page metadata; the date and authors cited
here are those printed on its first page.

For example, put $q=4$ and $\pi=2\zeta_5$.  Honda--Tate gives a simple
abelian surface over $\F_4$ with this Frobenius and
$\End^0(B_0)=\Q(\zeta_5)$.  Indeed the sole place above $2$ has
degree four, its slope is $1/2$, so its division-algebra invariant is
$4/2=0$ modulo $\Z$; all other invariants vanish.  The division algebra
is the field itself and the dimension is $4/2=2$.  The four CM types
on $\Q(\zeta_5)$ have trivial stabilizer, so all their reflex fields
are $\Q(\zeta_5)$.  Since $2$ has order four modulo five, the residue
field is $\F_{16}$, which cannot embed in $\F_4$.  Thus this isogeny
class, with its necessarily full CM endomorphism field, has no such
lifting.  This is the fully specified example in
\citet[Example (3.3)]{DOort2009}; the CM-type enumeration and the
polynomial $x^4+2x^3+4x^2+8x+16$ are independently checked in the
attached program.  Allowing a nonnormal domain changes the question:
\citet[Theorem (5.2)]{DOort2009} gives lifting up to isogeny without
extending the finite field in that broader setting.

\begin{lemma}[Proved: a power admits an exact-residue-field CM lifting]
\label{lem:d-power-cm-lift}
For every simple abelian variety $A_0/\F_q$ there exist $m\geq1$,
a number field $k$, a place $v$ of $k$ with $\kappa(v)=\F_q$, and a
CM abelian variety $B/k$, with its full CM field acting over $k$,
such that $B$ has good reduction at $v$ and
$B_v\sim_{\F_q}A_0^m$.  The integer $m$ is not required to equal one.
\end{lemma}

\begin{proof}
Write $D_0=\End^0(A_0)$, $F_0=\Q(\pi)$, and
$d^2=[D_0:F_0]$.  Honda--Tate gives
$2\dim A_0=d[F_0:\Q]$ and a maximal commutative CM subfield
$L_0\subset D_0$ with $[L_0:F_0]=d$.  These precise facts, including
the CM choice, are recalled at the beginning of
\citet[\S2, p.~7]{DCCO2008}.  Fix a Galois CM field $E_0$ containing
$L_0$ and also the imaginary quadratic field
$J=\Q(\sqrt{1-4p})$.  The polynomial $x^2-x+p$ has two distinct
roots modulo $p$, so $p$ splits in $J$.  Consequently complex
conjugation belongs to no decomposition group at $p$ in $E_0$.

Let $N$ be a common denominator of all slopes
$s_w=\ord_w(\pi)/\ord_w(q)$.  These belong to $[0,1]$ and satisfy
$s_{cw}=1-s_w$.  Start with the completion of $E_0$ at the fixed
$p$-adic place and enlarge it to a finite Galois extension $L_1/\Q_p$
containing $p^{1/N}$ and the $N$-th roots of unity.  Its ramification
index is divisible by $N$.  Let $f$ be its residue degree and choose
$\sigma\in\Gal(L_1/\Q_p)$ inducing arithmetic Frobenius on its
residue field.  Put $n=\operatorname{ord}(\sigma)$; then $f\mid n$.
If $U/\Q_p$ is unramified of degree $n$, then
$L=L_1U$ has the same ramification index and
\[
 \Gal(L/\Q_p)=I\rtimes\langle\tau\rangle,
 \qquad\operatorname{ord}(\tau)=n,
\]
where $\tau=(\sigma,\operatorname{Frob}_U)$ in the fibre product
over the cyclic residue group of $L_1$.  This proves the splitting
directly; it is also the splitting construction of
\citet[Remark 4.12 and Lemma 6.1]{DCCO2008}.

We need an enlargement which is still globally CM.  There exists a
totally real Galois number field $R$ whose completion at a specified
place is exactly $L$.  To prove this, take a primitive element of
the Galois extension $L/\Q_p$.  Approximate its minimal polynomial
$p$-adically and, simultaneously, a monic real polynomial with
distinct real roots, coefficient by coefficient by rational numbers.
For sufficiently close approximation Hensel's lemma supplies all
the roots inside $L$ and Krasner's lemma says that each root
generates $L$ over $\Q_p$.  At the real place all the roots remain
real.  The splitting field over $\Q$ is the required $R$.
Set $E=E_0R$.  It is Galois CM; its completion is $L$ and its
decomposition group is $\Delta=I\rtimes\langle\tau\rangle$.
Its inertia order $e=|I|$ is divisible by $N$, and $c\notin\Delta$.

Identify the embeddings of $E$ into $\overline\Q_p$ with
$G=\Gal(E/\Q)$.  The embeddings inducing one $p$-adic place form
a left coset $\Delta g$, of size $en$.  Such cosets occur in
disjoint conjugate pairs $\Delta g,c\Delta g$.  In $\Delta g$
choose exactly $e s_g$ orbits for the left action of
$\langle\tau\rangle$, each of size $n$.  This is possible since
$e s_g$ is an integer between zero and $e$.  On $c\Delta g$ choose
the complement of the conjugate choice.  The resulting subset
$\Phi\subset G$ is a CM type, is fixed by $\tau$, and satisfies
\[
 \frac{|\Phi\cap\Delta g|}{|\Delta|}=s_g.
\]
Thus it has exactly the required slopes.  If
$H=\operatorname{Stab}_G(\Phi)$, its reflex field is $E^H$.
Since $\tau\in H$, one has $I(\Delta\cap H)=\Delta$; the residue
degree of the induced place of $E^H$ is
$[\Delta:I(\Delta\cap H)]=1$.  Its residue field is therefore
$\F_p$, which embeds in the prescribed $\F_q$.

Finally $E$ splits $D_0$, since it contains $L_0$.  Put
$m=[E:F_0]/d$.  The splitting
$D_0\otimes_{F_0}E\simeq M_d(E)$ makes $E^d$ a right
$D_0$-module of rank $m$, with a commuting faithful left $E$-action.
It therefore embeds $E$ as a maximal commutative subfield of
$M_m(D_0)=\End^0(A_0^m)$.  The CM structure $(A_0^m,E,\Phi)$
satisfies both parts of the residual reflex condition.  Applying
\citet[Theorem 1.6, equivalently Proposition 5.4]{DCCO2008}
gives the asserted number-field CM lift with residue field
\emph{exactly} $\F_q$.
\end{proof}

\begin{corollary}[Proved: all Weil numbers occur, with variable CM data]
\label{cor:d-all-weil}
Let $\Gamma_q$ be the group of algebraic $p$-units whose complex
conjugates have absolute value $q^{w/2}$ for some integer $w$.
Every element of $\Gamma_q$ occurs as a Frobenius eigenvalue of the
good reduction of a tensor CM motive over some number field at a
place with residue field $\F_q$.  Every integral weight-one member
already occurs in $H^1$ of a CM abelian variety with such reduction.
No single number field, finite modulus, or primitive CM type is
asserted to work for all members.
\end{corollary}

\begin{proof}
For an integral weight-one Weil number, apply Honda--Tate and
Lemma~\ref{lem:d-power-cm-lift}.  For arbitrary $\alpha\in\Gamma_q$
of weight $w$, choose $a\geq0$ sufficiently large that
$\beta=\alpha q^a$ is integral and $b=w+2a>0$.  Any root
$\gamma=\beta^{1/b}$ is an algebraic integer, because it is
integral over the ring of algebraic integers, and every conjugate
of $\gamma$ has absolute value $q^{1/2}$.  It is therefore an
integral weight-one Weil number.  Choose the CM abelian variety
just constructed with $\gamma$ as a reduction eigenvalue.  Then
$(H^1(B))^{\otimes b}(a)$ has eigenvalue
$\gamma^bq^{-a}=\alpha$.  Conversely, tensor CM motives have Weil
eigenvalues with precisely these archimedean and non-$p$ integrality
properties, by good reduction and the Weil theorem for abelian
varieties, together with tensor products, duals and Tate twists.
\end{proof}

The earlier surface example also illustrates the power distinction
concretely.  For $q=4$, enlarge its CM field to
$E=\Q(\zeta_5,\sqrt{-3})$.  It embeds in
$M_2(\Q(\zeta_5))=\End^0(A_0^2)$.  Take the CM type induced from
$\Q(\sqrt{-3})$.  At $2$ its slope is $1/2$ at each place and its
reflex residue field is $\F_4$, so the residual reflex condition
holds.  Thus $A_0^2$ lifts up to $\F_4$-isogeny, although $A_0$
does not.  The corresponding finite-group calculation is included
in the attached verification program.

\subsection{D.2: toric identity components and the Taniyama group}
\label{subsec:d2}

\subsubsection{The exact general Hodge-class gate}

For either absolute-Hodge or motivated motives over a number field $k$,
write $\CM(k)$ for the Tannakian subcategory generated by Artin motives
and potentially CM abelian varieties.  Fix a complex embedding and the
Betti fibre functor.  The two categories are semisimple and the
motivic group of an object is reductive, possibly disconnected.

\begin{theorem}[Proved: D.2(a) for abelian type and its general equivalent]
\label{thm:d-toric-classification}
For a motive $M$ of abelian type, in either category,
\[
 G_{\mathrm{mot}}(M)^\circ\text{ is commutative}
 \quad\Longleftrightarrow\quad M\in\CM(k).
\]
For arbitrary motives the same classification is equivalent to the
following restricted Hodge-class assertion:

\emph{CM comparison case of the absolute-Hodge (respectively motivated
Hodge) conjecture.}  If $G_{\mathrm{mot}}(M)^\circ$ is commutative and
$N$ is a CM abelian motive over $\overline k$, every rational Hodge
morphism $H_B(M_{\overline k})\to H_B(N)$ is induced by an absolute
Hodge (respectively motivated) correspondence.

It is enough in this assertion to require one Hodge isomorphism to
some such $N$ to lift.  In the absolute-Hodge case this is a precise
restriction of \citet[Conjecture 8.24]{c_delignehodge1982}, rather
than a claim of equivalence with the unrestricted conjecture.
\end{theorem}

\begin{proof}
The Mumford--Tate group $\MT(H_B(M))$ is a connected subgroup of
$G_{\mathrm{mot}}(M)$, because every defining motivic tensor is a
Hodge tensor.  It is therefore contained in the identity component.
If that component is commutative, it is a torus and $H_B(M)$ is a
polarizable Hodge structure of CM type.  Every such Hodge structure
belongs to the tensor category generated by $H^1$ of CM abelian
varieties.  One way to see this explicitly is to use the character
lattice of $S$: the CM-type characters generate it integrally, as
proved below in the verification of R2.  Thus there exists a CM
abelian motive $N$ with a rational Hodge isomorphism
$H_B(M_{\overline k})\simeq H_B(N)$.

For abelian motives, Hodge realization is fully faithful by
\citet[Theorem 8.25, restating Main Theorem 2.11]{c_delignehodge1982}
in the absolute-Hodge category, and by
\citet[Theorem 0.6.2 and \S6.3]{c_andre1996} in the motivated
category.  The isomorphism therefore lifts to an isomorphism of
motives.  All the varieties and correspondences involved are defined
over some finite extension $k'/k$.  The adjunction between base
change and restriction of scalars has unit followed by trace equal
to $[k':k]$; dividing by this degree shows that $M$ is a direct summand
of $\operatorname{Ind}_{k'/k}(M_{k'})$.  Restriction of scalars of a
CM abelian variety is potentially CM, and induction of a Tate object
is Artin--Tate.  Hence this induced motive, and its summand $M$, lie
in $\CM(k)$.

Conversely, over $\overline k$ the category of CM abelian motives is
the category of representations of the Serre protorus.  Over $k$ it
is extended by the profinite Artin group.  Thus the identity component
of every object group in $\CM(k)$ is a quotient of a torus and is
commutative.

For general $M$ the same proof applies as soon as the comparison
isomorphism to $N$ lifts, and so the restricted conjecture implies
the classification.  If the classification holds, both $M$ and $N$
are of abelian type and the already cited full-faithfulness theorem
lifts every Hodge morphism between them.  This proves both directions
of the claimed equivalence, without assuming that a general geometric
motivic group is connected.  Such connectedness is itself conditional
in \citet[Propositions 8.22(a) and 8.23(b)]{c_delignehodge1982}.
\end{proof}

The 2026 theorem on Hodge classes of CM abelian varieties does not
settle this last comparison assertion: the classes here can lie on
$X\times B$ with $X$ an arbitrary smooth projective variety and $B$
CM abelian.  This distinction is reflected explicitly in
\citet[Remark 2.26(i), following Theorem 2.24]{DKahnOtsubo2026},
which treats the corresponding comparison condition for invertible
motives.  Thus D.2(a) is answered by the requested exact equivalence;
no proof of the unrestricted Hodge conjecture is being claimed.

\subsubsection{D.2(b): a non-metabelian Artin counterexample}

\begin{proposition}[Refuted: metabelian motives are all Taniyama quotients]
\label{prop:d-metabelian-false}
The category of Taniyama quotients is $\CM(\Q)$ and contains
motives with non-metabelian Galois group.  Its metabelian objects
are precisely the representations of the maximal metabelian quotient
of the Taniyama group.  Among motives of abelian type, this is
exactly the category with metabelian motivic Galois group.
\end{proposition}

\begin{proof}
The Taniyama group sits in
\begin{equation}
 1\longrightarrow S\longrightarrow\mathcal T
 \longrightarrow\Gal(\overline\Q/\Q)\longrightarrow1.
 \label{eq:d-taniyama}
\end{equation}
The quotient is the full absolute Galois group, not its abelianization.
Deligne's unnumbered main theorem on p.~262 identifies this group with
the motivic group of $\CM(\Q)$, preserving the Artin quotient, Hodge
cocharacter and finite-adelic splitting
\citep[pp.~261--264, main theorem and Lemmas 1--2]{c_delignetaniyama1982}.

Let $f(x)=x^4-x-1$, $\alpha$ a root and
$M=h^0(\Spec\Q(\alpha))$.  Reduction modulo two gives the irreducible
$x^4+x+1$: it has no linear factor and is not divisible by the only
irreducible quadratic $x^2+x+1$.  Thus $f$ is irreducible over $\Q$.
Its discriminant is $-283$.  Modulo seven it factors as
\[
 (x-3)(x^3+3x^2+2x-2),
\]
with the cubic irreducible.  The splitting-field group contains a
four-cycle and a three-cycle.  Its order is divisible by twelve,
and the four-cycle is odd, so the subgroup is $S_4$, not $A_4$.
The permutation representation on the four embeddings is faithful.
Consequently $G_{\mathrm{mot}}(M)=S_4$ in both categories, by
\citet[Proposition 8.17]{c_delignehodge1982} and the same Artin
realization in motivated motives.  But
\[
 [S_4,S_4]=A_4,\qquad[A_4,A_4]=V_4\ne1.
\]
All these finite calculations are checked by the attached program.
This is an Artin motive and hence a quotient representation of
$\mathcal T$, but its motivic group is not metabelian.  Its identity
component is trivial, so it also refutes the inference ``commutative
identity component implies metabelian'' in the supplied reading of R4.

The exact replacement is the category $\CM(\Q)=\Rep_\Q(\mathcal T)$.
Inside it, the objects with metabelian group are exactly
\[
 \Rep_\Q\!\left(\mathcal T/
       \overline{[[\mathcal T,\mathcal T],[\mathcal T,\mathcal T]]}\right),
\]
where the bar means the closed normal subgroup generated by the
indicated commutators, equivalently the intersection of the kernels
of all metabelian algebraic quotients.  This equality follows directly
by factoring a representation through its scheme-theoretic image.
For motives of abelian type, a metabelian reductive connected group
is a torus, so Theorem~\ref{thm:d-toric-classification} places them
inside this category.  For arbitrary motives, the additional step is
exactly the comparison gate in that theorem.
\end{proof}

\subsubsection{D.2(c): comparison of the Hodge structures at two places}

\begin{proposition}[Known and rechecked: the Taniyama comparison]
\label{prop:d-hodge-comparison}
The two Betti fibre functors on geometric CM motives at
$\widetilde\xi$ and $\widetilde\xi'=\sigma\widetilde\xi$
are related by the fibre of the Taniyama group above
$\sigma|_{\overline\Q}$.  Their Hodge types change by the
$\Aut(\C)$ action and their finite-adelic comparison is supplied
by its canonical adelic section.  Formula~\eqref{eq:d-conjugate-cm-type}
gives the explicit CM-type and polarization change.
\end{proposition}

\begin{proof}
Let $\widetilde\xi:k\hookrightarrow\C$ and
$\widetilde\xi'=\sigma\widetilde\xi$, where
$\sigma\in\Aut(\C)$, and let $\tau$ be its restriction to
$\overline\Q$.  The fibre $\mathcal T_\tau$ of
\eqref{eq:d-taniyama} is an $S$-torsor.  It is the tensor-isomorphism
torsor between the Betti fibre functors at the two embeddings on
geometric CM motives.  This identification is
\citet[Proposition 8.23(c)--(d)]{c_delignehodge1982} specialized by
Deligne's Taniyama theorem.  A tensor trivialization transports the
underlying rational vector spaces, and the Hodge cocharacter is
transported to its $\sigma$-conjugate.  Such a rational trivialization
need not be canonical, or even exist; the adelic splitting supplies
the canonical finite-adelic comparison.  In particular, the answer
does not identify the two rational Hodge structures without changing
the CM type.

More explicitly, use homology and write a polarized CM abelian
variety at $\widetilde\xi$ in the form
$(E,\Phi;\mathfrak a,t)$.  For
$f\in f_\Phi(\sigma)\subset\mathbb A_{f,E}^{\times}/E^\times$,
its conjugate has type
\begin{equation}
 (E,\sigma\Phi;
       f\mathfrak a,\ t\chi_{\mathrm{cyc}}(\sigma)/(f\overline f)).
 \label{eq:d-conjugate-cm-type}
\end{equation}
There is an $E$-linear rational homology isomorphism $\alpha$ for
which $(\alpha\otimes1)f$ is the finite-adelic comparison, with
the displayed polarization factor.  These are
\citet[Theorems 1.1 and 1.3]{DMilnePedestrians}.
Applying tensor products, duals and projectors gives the comparison
for every CM motive.  On cohomology one dualizes the formula; this
accounts for the sign of the norm weight in 15:165.

At a real embedding the compatible de Rham structure is the fixed
real form of the Betti comparison under
$F_\infty\otimes$ complex conjugation.  This verifies the real-place
condition used from R4: the involution acts on the Betti factor and
coefficient conjugation acts on $\C$, so their fixed subspace is
the real de Rham realization.  It does not remove the torsor ambiguity
of arbitrary rational fibre functors examined next.
\end{proof}

\subsection{D.3: the finite-place kernel for the actual Serre torus}
\label{subsec:d3}

Let $K/\Q$ be a finite Galois CM extension, $F=K^+$, and
$r=[F:\Q]$.  Set
\[
 \Sha^1_f(\Q,S^K)=
 \ker\!\left(H^1(\Q,S^K)\longrightarrow
                      \prod_{\ell<\infty}H^1(\Q_\ell,S^K)\right).
\]
The subscript $f$ is essential: this is not the ordinary
Tate--Shafarevich group, which also tests the real place.

\begin{theorem}[Proved: finite-place kernel and fibre functors]
\label{thm:d-serre-finite-kernel}
There is a canonical identification
\[
 \Sha^1_f(\Q,S^K)
 \simeq
 \left\{(\epsilon_v)_{v\mid\infty}\in\F_2^r:
                            \sum_v\epsilon_v=0\right\}.
\]
For the geometric neutral category of CM motives whose CM fields
are subfields of $K$, with Betti functor as origin, this group
classifies the isomorphism classes of rational fibre functors that
become isomorphic to Betti at every finite completion.  It has
$2^{r-1}$ elements.  Thus finite comparisons single out Betti only
when $r=1$.  Adding a real comparison singles it out for every $r$.
\end{theorem}

\begin{proof}
There is an exact sequence of tori
\begin{equation}
 1\longrightarrow\Res_{F/\Q}\Gm
 \xrightarrow{a\mapsto(a,N_{F/\Q}(a)^{-1})}
 \Res_{K/\Q}\Gm\times\Gm
 \longrightarrow S^K\longrightarrow1.
 \label{eq:d-serre-resolution}
\end{equation}
Its character sequence says exactly that the coefficients $f_\sigma$
have constant pair sums $f_\sigma+f_{c\sigma}$.  Hilbert 90 gives
\begin{equation}
 H^1(\Q,S^K)=
 \ker\left(\Br(F)\xrightarrow{(\res,-\cor)}
                         \Br(K)\oplus\Br(\Q)\right).
 \label{eq:d-brauer-kernel}
\end{equation}
The same calculation holds over every completion.  This is
\citet[Equation (23) and Lemma 3.1]{AMilne2003}, with the sign
of the second map arising from the inverse norm.

An element in the finite-place kernel corresponds under
\eqref{eq:d-brauer-kernel} to a Brauer class $b$ whose finite local
invariants vanish.  Its real invariants are
$\operatorname{inv}_v(b)=\epsilon_v/2$, with
$\epsilon_v\in\{0,1\}$.  Global Brauer reciprocity imposes
$\sum_v\epsilon_v=0$ in $\F_2$.  Conversely, every such even
pattern is realized by a unique global Brauer class.  Its restriction
to $K$ vanishes because the real places become complex and the
finite invariants are zero.  Its corestriction to $\Q$ also vanishes:
its only possible nonzero invariant is at infinity and is the sum
of the displayed real invariants.  This proves both existence and
uniqueness in the formula.  It also proves that restriction to
$H^1(\R,S^K)$ is an isomorphism on this kernel and that the kernel
over \emph{all} places is zero, in agreement with
\citet[Proposition 3.2]{AMilne2003}.

The geometric CM category under discussion has Tannaka group $S^K$:
all its Hodge characters satisfy the pair-sum condition, and CM
types generate that lattice integrally.  The absolute-Hodge and
motivated versions agree here by the full-faithfulness results cited
in Theorem~\ref{thm:d-toric-classification}.  The usual torsor
classification sends a fibre functor $\omega$ to
$\Isom^\otimes(\omega_B,\omega)$ in $H^1(\Q,S^K)$; conversely an
$S^K$-torsor twists $\omega_B$.  The local comparison requirement
is exactly the displayed kernel.  Since $S^K$ is commutative these
isomorphism classes form a group after choosing Betti as origin,
or intrinsically a torsor under that group.
\end{proof}

\subsubsection*{Verification of R3 and the missing comparison in Lead D3}
Put $U=\Res_{F/\Q}(\Res^1_{K/F}\Gm)$.  There is an exact sequence
\[
 1\longrightarrow\Gm\longrightarrow S^K
       \xrightarrow{(x,z)\mapsto x/\overline x}U\longrightarrow1.
\]
Thus $H^1(\Q,S^K)$ injects into
$H^1(\Q,U)=F^\times/N_{K/F}K^\times$.  The connecting map to
$\Br(\Q)$ is the corestriction of the cyclic algebra
$(K/F,a)$.  A norm class locally trivial at every finite place has
only real invariants; reciprocity imposes the same even-sign
condition, and then its corestriction vanishes.  Therefore this
injection induces an isomorphism on the finite-place kernels.
This derives R3 by the cyclic norm/Brauer exact sequence and proves
the additional step required to use it for $S^K$.  Identifying the
two tori themselves would be incorrect.

\begin{example}[Verified: a non-Betti form with all finite comparisons]
\label{ex:d-non-betti}
Take $F=\Q(\sqrt5)$ and $K=F(i)$.  The quaternion algebra
$(-1,-1)_F$ is nonsplit at both real embeddings.  It is split at
every finite place: over $\Q$ Hamilton's algebra is ramified only
at $2$ and infinity, and $2$ is inert in $F$ because
$x^2-x-1$ is irreducible modulo two.  The degree-two local extension
at $2$ multiplies the invariant $1/2$ by two.  The algebra is split
by $K$, and its corestriction to $\Q$ is trivial since its two real
invariants sum to zero in $\Q/\Z$.  It is the nonzero element of
$\Sha^1_f(\Q,S^K)\simeq\Z/2$.  Equivalently, the corresponding
$U$-torsor is $x^2+y^2=-1$ over $F$.  Twisting Betti by its lift to
an $S^K$-torsor supplies the asserted different rational fibre
functor.  The irreducibility calculation is recorded in the program.
\end{example}

The word ``geometric'' in the theorem cannot be omitted.  Over a
number field the arithmetic CM category also contains the Artin
quotient in \eqref{eq:d-taniyama}.  Its fibre functors are torsors
under that full arithmetic group, or under the particular quotient
generated by the chosen motives; replacing it by $S^K$ without
fixing the Artin data changes the question.  Likewise a smaller CM
subcategory has its actual quotient torus as band.  The computation
above addresses the full geometric category and finite $K$, not
an unexamined inverse limit over all CM fields.

Finally, this is an ambiguity in $H^1$ between fibre functors of an
already neutral category.  The $\mu_2$ at infinity in 15:11 is an
$H^2$ obstruction to neutralizing a different finite-field gerbe.
The former category has the Betti real fibre functor, while the
latter may have no ordinary real fibre functor at all.  Their common
archimedean origin does not identify their cohomology classes.  The
same distinction must be retained when comparing the real functors
of SA06-10: once a real functor exists its forms are $H^1$-classes;
existence itself is the gerbe's $H^2$ question.

\subsection{D.4: the image of the type norm, including its integral lattice}
\label{subsec:d4}

Let $I_E=\Hom(E,\overline\Q)$ and
$I_K=\Hom(K,\overline\Q)$.  Let $\Sigma\subset I_E\times I_K$
be the Galois-invariant CM structure in 15:147--151, and use the
convention in which its type norm has exponent matrix
\[
 A=(a_{\sigma,\tau})_{\sigma\in I_E,\tau\in I_K},
 \qquad a_{\sigma,\tau}=1_\Sigma(\sigma,\tau).
\]
Inverting all exponents to switch between homology and cohomology
does not change the image.  The source map $T_K\to T_E$ factors
through $S_K^0$ as in R2.

\begin{theorem}[Proved: the complete lattice and rank answer to 15:151]
\label{thm:d-reflex-image}
The image $I_\Sigma$ of $\psi_\Sigma$ is the connected $\Q$-torus
with character and cocharacter lattices
\begin{align*}
 X^*(I_\Sigma)&=\Z^{I_E}/\ker(A^{\mathsf t})
                  \simeq\im(A^{\mathsf t})\subset\Z^{I_K},\\
 X_*(I_\Sigma)&=(\im A\otimes\Q)\cap\Z^{I_E}.
\end{align*}
In particular $\dim I_\Sigma=\operatorname{rank}_\Q(A)$.
If $E$ is CM of degree $2g$, put
\[
 C_E=\{x\in\Res_{E/\Q}\Gm:
              x\overline x\in\Gm\subset\Res_{E^+/\Q}\Gm\}.
\]
Then $I_\Sigma\subset C_E$, and the following conditions are equivalent:
\begin{enumerate}[label=(\roman*)]
 \item $I_\Sigma=C_E$;
 \item $\operatorname{rank}_\Q(A)=g+1$;
 \item every rational row relation of $A$ is generated by
 $e_\sigma+e_{c\sigma}-e_\tau-e_{c\tau}$.
\end{enumerate}
These are effective necessary and sufficient conditions for every CM
type, including non-Galois and imprimitive types.  For the CM type
$\Phi$ represented by $\Sigma$, $I_\Sigma=\MT(E,\Phi)$.
\end{theorem}

\begin{proof}
Over $\overline\Q$, the map is
$y_\sigma=\prod_\tau x_\tau^{a_{\sigma,\tau}}$, so its character map
is $A^{\mathsf t}$.  The characters vanishing on its image are
exactly $\ker A^{\mathsf t}$.  That kernel is saturated, because the
target $\Z^{I_K}$ is torsion-free.  Hence the displayed quotient
is a free lattice and describes a connected torus.  Dually, the
cocharacters of a subtorus are the saturation of the image on
cocharacters, which gives the second formula.  Omitting this
saturation can give an incorrect integral lattice even though it
preserves the rank.

For a CM structure each conjugate row pair has sum the all-ones
row.  Therefore every vector
$e_\sigma+e_{c\sigma}-e_\tau-e_{c\tau}$ lies in the character
kernel.  These are exactly the defining relations of $C_E$, whose
dimension is $g+1$.  Inclusion follows.  A connected closed subtorus
of the same dimension is the whole torus, proving (i)$\Leftrightarrow$(ii).
The relation space just described has dimension $g-1$; equality of
its dimension with the full row-kernel dimension proves (iii).
Over $\Z$ the exact answer remains the displayed kernel, computable
by Smith normal form, rather than just a choice of rational basis.

The Hodge cocharacter $\mu_\Phi$ of the rank-one $E$-module has
coordinates $1_\Phi$.  The columns of the reflex-norm matrix are its
Galois conjugates, with repetitions if the chosen source contains
the reflex field properly.  Hence the image is the smallest rational
subtorus containing $\mu_\Phi$.  A rational algebraic subgroup
containing the Hodge homomorphism contains $\mu_\Phi$ and every
Galois conjugate; conversely the torus they generate contains both
$\mu_\Phi$ and its complex conjugate and therefore contains the
Hodge homomorphism.  This is the defining minimality property of
the Mumford--Tate group.  It proves the equality without assuming
nondegeneracy.  The reflex norm identity is stated in
\citet[Proposition 1.23 and \S4, p.~35]{DMilneCM2020}; its
universal cocharacter property is Proposition 4.19, and the CM
criterion for the Mumford--Tate group is Corollary 4.7 there.
\end{proof}

\begin{corollary}[Proved: the compact part in the manuscript's reduction]
\label{cor:d-compact-image}
Let $V'_K=\ker(N_{K/\Q}:S^K\to\Gm)$, and let
$\nu:C_E\to\Gm$ send $x$ to its common product
$x_\sigma x_{c\sigma}$.  Then
\[
 \psi_\Sigma(V'_K)=I_\Sigma\cap\ker\nu,\qquad
 \dim\psi_\Sigma(V'_K)=\operatorname{rank}_\Q A-1,
\]
and the character lattice of this image is
\[
 \im(A^{\mathsf t})/\Z\mathbf1_{I_K}.
\]
The image is all of the norm-one torus of $E/E^+$ precisely when
$\operatorname{rank}_\Q A=g+1$.
\end{corollary}

\begin{proof}
The pair-sum identity gives $\nu\circ\psi_\Sigma=N_{K/\Q}$.
If an element of $I_\Sigma$ has $\nu=1$, any fppf-local lift to
$S^K$ consequently has norm one.  This proves equality of the
fppf images.  The pullback of $\nu$ is the all-ones vector,
which is primitive even in $\Z^{I_K}$, hence also in
$\im(A^{\mathsf t})$.  Its kernel in $I_\Sigma$ is therefore
connected, and the exact character sequence gives the formula
and rank.  The norm-one torus of $E/E^+$ has dimension $g$;
equality follows from the same dimension criterion as above.
Over $\R$ it is compact, which explains the manuscript's reduction
to this part.  The weight cocharacter supplies the omitted scalar
direction; its intersection with the compact part is $\mu_2$,
so this decomposition is an isogeny, not an asserted integral
direct-product decomposition.
\end{proof}

\paragraph{A complete Fourier test in the abelian Galois case.}
If $E/\Q$ is abelian Galois with group $G$ and central conjugation
$c$, use all translates of $1_\Phi$ as columns.  Diagonalizing the
convolution matrix by the characters of $G$ gives
\[
 \dim\MT(E,\Phi)=1+
 \#\left\{\chi\in\widehat G:\chi(c)=-1,
                         \ \sum_{\phi\in\Phi}\chi(\phi)\ne0\right\}.
\]
Indeed the trivial Fourier coefficient is $g$; every nontrivial
even character has sum zero on $\Phi$, since doubling that sum
gives its sum on $G$; and the remaining eigenvalues are precisely
the displayed odd sums.  Thus the full norm-compatible torus
occurs exactly when every odd character has nonzero type sum.
This criterion, together with Theorem~\ref{thm:d-reflex-image}, is
a classification test for all cases, rather than a list of only
the nondegenerate examples.

For example the imaginary quadratic type has rank two, the cyclic
quartic type $\{0,1\}\subset\Z/4$ has rank three, the cyclic sextic
type $\{0,1,2\}\subset\Z/6$ has rank four, and the induced sextic
type $\{0,2,4\}$ has rank two.  For a biquadratic quartic field every
type is induced from an imaginary quadratic subfield and has rank
two, less than $\dim C_E=3$.  The attached program prints all the
matrices, exact ranks, kernels and Smith invariants in these cases.
If $E$ itself is not CM but the CM structure is pulled back from
a minimal CM subfield $E_0$, the same matrix formula applies and
the image is the diagonal image of $\MT(E_0,\Phi_0)$ in $T_E$.
The notation $C_E$ above is used only when $E$ is CM.

\paragraph{Verification of the portions of R2 used here.}
For a CM field with $r$ conjugate pairs, the character lattice of
the connected Serre quotient is
\citep[\S1, equations (1.1)--(1.3)]{DMilneShih1982}
\[
 L_S=\{(n_\sigma): n_\sigma+n_{c\sigma}
                            \text{ is independent of }\sigma\}.
\]
Choose one CM type $\Phi$.  A basis of this lattice is $1_\Phi$
together with the $r$ vectors $e_\sigma-e_{c\sigma}$, one for each
pair with $\sigma\in\Phi$.  Each latter vector is the difference
between two CM-type indicators obtained by flipping one pair.
Consequently all CM types generate $L_S$ over $\Z$, not merely
over $\Q$.  This proves the simultaneous faithfulness invoked in
R2 and in Theorem~\ref{thm:d-serre-finite-kernel}.  The same pair-sum
identity proves that $\psi_\Sigma$ annihilates the norm-one subtorus
from the maximal real subfield and factors through $S_K^0$.
The Galois orbit of a CM type is a subset of the $2^r$ possible
choices, so its reflex field has degree at most $2^r$.  Finally the
character argument for $\MT$ above proves the part of R2 concerning
the universal Hodge cocharacter; it does not identify a fixed
$T_{\mathfrak m}$ with $S^K$ in spite of the finite unit quotient.
